%============================================================================
% NATURE OF THIS COMPANION
%============================================================================

\noindent\textit{\textbf{On the nature of this companion.}
This document is the executive summary of a voluntary contribution
produced by members of the academic community acting in their
scientific and technical capacity, addressed to the
institution's governing bodies (the unit-level deliberative body,
the institution-wide council, and \loc{governing-board}) for
consideration and, if found substantive, propagation by
\loc{governing-board} to all concerned components.}

\bigskip

%============================================================================
% KEY MESSAGES BOX
%============================================================================

\begin{tcolorbox}[
    colback=NavyBlue!6!white,
    colframe=NavyBlue!80!black,
    boxrule=0.8pt,
    arc=2pt,
    left=10pt, right=10pt, top=8pt, bottom=8pt,
    title={\textbf{Key messages}}
]
\begin{itemize}[nosep, leftmargin=1em]
    \item \textbf{Digital freedom is the rule} in a university; restrictions
    must be proportionate, documented, and limited to what is strictly
    necessary. When a contract or a regulation would invert that order,
    the institution's first reflex is to renegotiate or to engage in
    reform --- not to comply by default (the main report's
    liberty-first rule).
    \item \textbf{Segment to protect, liberate to excel:} a multi-zone
    architecture with Trusted Activity Contexts strengthens security
    for sensitive data while granting autonomy to research and teaching.
    \item \textbf{Democratic governance of digital policy:} a Digital
    Policy Commission chaired by an elected academic formulates proposals
    endorsed by the unit-level deliberative bodies and the
    institution-wide council.
    \item \textbf{Proven and progressive:} the approach is validated by
    peer research institutions, international standards, and the
    applicable regulatory framework; it can be deployed incrementally, starting with a pilot
    on selected entities.
\end{itemize}
\end{tcolorbox}

\bigskip

%============================================================================
\section{The university spirit: openness as a foundational principle}
%============================================================================

From the \emph{Magna Charta Universitatum} (Bologna, 1988) to the UNESCO
\emph{Recommendation on Science and Scientific Researchers} (2017), a
universal consensus holds: the university is a space of \textbf{intellectual
freedom}, \textbf{free movement of people and ideas}, and
\textbf{autonomy} for academic staff. This openness is not a luxury ---
it is the very condition under which knowledge is produced and transmitted.
Today, the digital space is the primary space where this movement occurs.
\textbf{A digital policy that is restrictive by default
is contrary to the university spirit.}

\begin{principlebox}[Founding principle]
A university requires administration, but is not a corporation;
university governance must optimise for freedom for ideas and persons,
physical or digital. The digital infrastructure should reflect this
fundamental distinction.

\smallskip
In a university, \textbf{digital freedom is the rule}. Any restriction must
be \emph{explicitly justified} by a documented risk, \emph{proportionate}
to that risk, and \emph{limited to the strictly necessary scope} --- a
specific laboratory, a specific project, a specific dataset --- never
applied uniformly across the entire institution.

\smallskip
This principle reflects the risk-based, proportionate-measures requirement
that runs through modern data-protection and cybersecurity frameworks --- in
the European Union, \gls{gdpr} Article~32 (\emph{measures appropriate to the
risk}) and the NIST Cybersecurity Framework~2.0 (\emph{risk-based approach})
are canonical references --- and it is consistent with the practice of leading
research universities.
\end{principlebox}

\mainref{the sections \nameref{sec:stakeholders}, \nameref{sec:introduction}
and \nameref{sec:specificites}}

%============================================================================
\section{What the university spirit demands of digital infrastructure}
%============================================================================

\begin{table}[h!]
\centering\small
\newcommand{\ok}{\textcolor{OliveGreen}{\ding{51}}\,}%
\newcommand{\warn}{\textcolor{orange}{\ding{72}}\,}%
\newcommand{\nok}{\textcolor{BrickRed}{\ding{55}}\,}%
\begin{tabularx}{\textwidth}{>{\bfseries\raggedright\arraybackslash}X X}
\toprule
\textbf{Historical university principle} &
  \textbf{Digital translation (current status --- keep one marker per element)} \\
\midrule
Free movement of persons &
  \ok\warn\nok Functional BYOD,
  \ok\warn\nok eduroam,
  \ok\warn\nok seamless researcher mobility \\
\addlinespace
Free movement of ideas &
  \ok\warn\nok unfiltered Internet access,
  \ok\warn\nok major cloud-based AI services,
  \ok\warn\nok worldwide scientific sources \\
\addlinespace
Teacher-researcher autonomy &
  \ok\warn\nok local admin privileges (root/admin),
  \ok\warn\nok freedom of tool choice,
  \ok\warn\nok configurable environments \\
\addlinespace
Scientific openness (Open Science) &
  \ok\warn\nok access to data platforms,
  \ok\warn\nok code repositories,
  \ok\warn\nok open-access publications \\
\addlinespace
Realistic training environments &
  \ok\warn\nok configurable labs,
  \ok\warn\nok cyber-ranges,
  \ok\warn\nok API access for AI courses \\
\bottomrule
\end{tabularx}

\smallskip
\footnotesize
\ok Available \qquad
\warn To be reinforced \qquad
\nok Not available
\par\smallskip
\textit{This table is generic: each element carries all three markers. The
adopting institution keeps, for every element, the single marker matching its
own reality and deletes the other two.}
\normalsize
\end{table}

%============================================================================
\section{Effects of uniform restriction policies}
%============================================================================

A policy that treats financial data and research GPU clusters
identically generates several documented effects: it concentrates
risk on a single perimeter (a breach can affect the institution
as a whole --- as recent breaches in the sector
have shown); it correlates with the
emergence of Shadow IT (researchers may turn to unsupervised
solutions); it is associated with talent mobility patterns toward
institutions with more flexible IT environments; and it affects
competitiveness in European and international grant calls.
Restrictive policies can also expose the institution to
\emph{litigation based on fundamental rights} (academic freedom,
freedom of expression and information, dignity at work).

\begin{principlebox}[Guiding principle for institutional resilience]
A university information system must not attempt to prevent all errors,
experiments, or unforeseen uses; it must \textbf{prevent such events from
compromising the institution as a whole}.

\smallskip
\textbf{In short:} local freedom, incident confinement, protection of the
core, global resilience.
\end{principlebox}

\begin{alertbox}[What the most publicised recent university incidents actually show]
The most publicised university breaches of recent years --- whether at
large national research universities or through shared platform and
enterprise-software vulnerabilities such as those affecting widely used
managed file-transfer and ERP products --- share \textbf{the same root causes}:
third-party/supply-chain vulnerabilities, credential compromise, and
centralised perimeters without internal boundaries.
\textbf{None of the major recent data breaches analysed here were caused by research-zone
autonomy.} Known incidents involving research infrastructure --- such as
the 2020 European supercomputer attacks, where stolen SSH credentials
were used for cryptomining on \gls{hpc} systems at several national
academic computing centres~\cite{eu-supercomputers-2020}~--- involved unauthorized access rather than data exfiltration, and
would have been mitigated by the Trusted Activity Context mechanisms
proposed here: credentials scoped to a specific activity context cannot
grant lateral access to other contexts or to the infrastructure core.
They demonstrate the urgent need for \emph{segmentation and
context-bounded access} --- precisely what this proposal advocates.
\end{alertbox}

\mainref{the sections \nameref{subsec:shadow-it}, \nameref{subsec:dependability},
\nameref{sec:defaillances-mediatiques} and \nameref{sec:couts}}

%============================================================================
\section{The proposal: segment to liberate, Trusted Activity Contexts for individualised protection}
%============================================================================

The logic is simple: \textbf{secure strongly what must be secured}
(administrative, financial, HR data) \textbf{to liberate everything else}.
This is what research universities and national laboratories
do --- and, through their national research-and-education
networks, what institutions of far more modest scale do as well ---
and what NIST, ISO~27001 and \gls{cobit} recommend. These are not a
flagship's choices to borrow but a minimal set every institution
offering a digital infrastructure must discharge; only the
instantiation varies with size, not the obligation.

The proposed architecture rests on two complementary mechanisms:
\emph{segmentation} controls where incidents can spread (network
boundaries); \textbf{Trusted Activity Contexts} provide each user with
a validated environment --- assembling data, tools, services, and access
rights --- adapted to their current role and task. Together, they ensure
that a breach in one context cannot cascade to the rest of the
institution.

Conceptually, the university information system is organised in
\textbf{concentric layers} (the executive view of the five structural
layers detailed in the main report): a small, strongly protected
\textbf{infrastructure core} (identity, network backbone, backups,
supervision) that forms the \emph{dependability foundation} of the
institution; \textbf{common services} (messaging, \gls{lms}, collaboration)
available to all; \textbf{autonomous zones} (research, teaching) with
local freedom and bounded impact; \textbf{exposed zones} (public-facing
services and portals, isolated from the internal system by DMZ-type
architectures); and \textbf{user networks} (\gls{byod},
Wi-Fi) that are open but isolated from the core. Protecting the small
critical core is what makes it possible to liberate everything else.

The six principal operational zones below map onto these layers
(a seventh IoT/Building zone is detailed in the technical report):

\begin{table}[h!]
\centering\small
\renewcommand{\arraystretch}{1.15}
\begin{tabularx}{\textwidth}{
  >{\bfseries\raggedright\arraybackslash}p{2.9cm}
  >{\raggedright\arraybackslash}X
  >{\raggedright\arraybackslash}p{4.2cm}}
\toprule
\textbf{Zone} & \textbf{Philosophy} & \textbf{Constraint level} \\
\midrule
Research &
  \textbf{Maximum openness} --- autonomy, root access, worldwide AI cloud,
  local server management &
  Minimal: responsibility charter, monitoring \\
\addlinespace
Teaching &
  \textbf{Broad openness} --- discipline-specific environments
  for all faculties, AI APIs, cyber-ranges, VDI &
  Standard: pedagogical framework \\
\addlinespace
Staff BYOD &
  \textbf{Free Internet} + progressive access to institutional resources &
  \gls{mfa} + minimal posture check \\
\addlinespace
Research-data path (e.g.\ Science DMZ) &
  \textbf{Maximum performance} --- high-throughput transfers without
  stateful firewall &
  \gls{acl} + dedicated \gls{ids} \\
\addlinespace
Administrative &
  \textbf{Maximum protection} --- this is where, \emph{and only where},
  restrictions are concentrated &
  Strict: firewall, \gls{ids}/\gls{ips}, proxy, full \gls{gdpr} \\
\addlinespace
Sensitive enclave &
  Case by case, \textbf{at project level} &
  Only on explicit, documented justification \\
\bottomrule
\end{tabularx}
\end{table}

\noindent\textbf{Key point:} restrictions at laboratory or project level apply
\textbf{only when explicit and documented reasons require them} (medical data,
classified projects, specific funding constraints such as NIST~SP~800-171). In
the absence of such justification, the laboratory benefits from the \textbf{maximum
openness policy} of the research zone. The default is freedom; the exception is
restriction.

\smallskip
\noindent\emph{All zones lie within the cybersecurity-baseline perimeter
applicable to the institution under its entity classification; calibration of
controls follows the risk-proportionality principle of that regime --- the
\gls{nis2} Article~21 measure families in the worked EU instantiation --- not
scope reduction. The function this discharges is the cybersecurity
risk-governance obligation (F2) of the Regulatory Applicability annex
(Part~\ref{part:reg-applicability} of the complete dossier); read your own regime's instrument for
that obligation in the regime typology there.}

\mainref{the sections \nameref{sec:etat-art}, \nameref{sec:exemples},
\nameref{sec:ia} and \nameref{sec:propositions-infra}}

%============================================================================
%\newpage
\section{Governance aligned with the university spirit}
%============================================================================

\begin{itemize}
  \item \textbf{Digital Policy Commission} chaired by an \emph{elected}
    academic, composed of representatives from the constituent academic
    units (departments and interdisciplinary centres) across the
    institution. The Commission formulates policy proposals that are
    submitted to the unit-level deliberative bodies and then to the
    institution-wide council for endorsement. The IT-management lead
    and the information-security-management lead submit proposals; the
    Commission decides --- aligned with \gls{cobit}~2019's separation of governance
    (strategic direction) from management (operational execution)
    and with established academic-governance practice;
  \item \textbf{Responsibility charters} (as practised across research
    institutions): autonomy
    is formalised and accountable, replacing informal exceptions;
  \item \textbf{Contextual role management:} a single individual
    (e.g.\ a researcher who also teaches and administers a laboratory)
    activates only the relevant role for each task --- never carrying all
    privileges simultaneously. This enforces least privilege, improves
    auditability, and maps naturally onto the zone structure;
  \item \textbf{Open-source and open-science tools as default
    institutional choices}, reducing vendor lock-in, increasing
    transparency, and aligning with public-sector open-source
    strategies (e.g., the European Commission's in the EU);
  \item \textbf{Independent Commission for Reasonable Adjustments
    of Digital Policies} (\gls{icradp}): any employee can
    challenge a disproportionate restriction --- for medical, personal,
    professional reasons, or infringement of fundamental freedoms ---
    in accordance with the EU Charter of Fundamental Rights, the
    CRPD Convention and the case law of the \gls{cjeu} and the \gls{ecthr}
    on the proportionality of digital restrictions. Where the institution
    already operates a reasonable-adjustments mechanism (supporting
    students and staff), the \gls{icradp} extends that principle to
    digital policies; where none exists, it establishes the principle
    directly. Refusals must be justified in writing.
\end{itemize}

\medskip
\noindent\emph{Continuous-improvement mechanism.} Together, the
Digital Policy Commission, the unit-level deliberative bodies, the
responsibility charters and the ICRADP constitute the institution's
continuous-improvement mechanism, materialised in the annual
compliance and maturity report submitted to the competent cybersecurity
authority~\cite{ilr-ul-essential} pursuant to the continuous-improvement
provision of the applicable regime --- see the Regulatory Applicability annex
(Part~\ref{part:reg-applicability} of the complete dossier), the accountability-and-conformity
obligation (F6) under your regime profile, instantiated in the worked EU
example by the \loc{nis2-national-transposition}~\cite{pl-8364}.

\mainref{the sections \nameref{subsec:gouvernance-federee}, \nameref{subsubsec:contextual-roles},
\nameref{sec:propositions} and \nameref{sec:commission-ajustements}}

%============================================================================
\section{Feasibility: progressive, controlled, affordable}
%============================================================================

\begin{table}[h!]
\centering\small
\begin{tabularx}{\textwidth}{
  >{\bfseries}l X}
\toprule
\textbf{Step} & \textbf{Content} \\
\midrule
1. Governance &
  Digital Policy Commission, responsibility charters, service catalogue,
  formalised exception path \\
\addlinespace
2. Pilot &
  Segmented architecture and Trusted Activity Contexts deployed on
  a pilot technical department, a pilot research centre, and the
  \loc{pilot-entities} ---
  high-throughput research-data path (Science DMZ), \gls{byod} zone, open-source tools \\
\addlinespace
3. Progressive extension &
  Gradual rollout to other departments, faculties and centres,
  progressive \gls{zt}, AI research zone,
  audit and continuous improvement \\
\bottomrule
\end{tabularx}
\end{table}

\noindent
No ``big bang'' required. The governance step is organisational and
achievable with existing staff. The pilot and progressive extension
rely on the \loc{nren-name} and its international
research-backbone services, open-source tools,
and platforms offered free of charge for academic research.

\medskip
\noindent\textbf{Scale-appropriate reading.} Where public documentation
evidences no Science DMZ --- nor any equivalent research-and-education-network-operated
research-data path --- deployed, neither at the institution nor at the
\loc{nren-name} level, the substrate is nonetheless frequently
already operational: high-capacity backbone connectivity to the
international research backbone through the \loc{nren-name}, an \gls{hpc} facility moving
petabyte-scale traffic, and \loc{nren-csirt} for incident
response. For institutions of modest scale, the pragmatic form is
therefore not an institution-only deployment of the Science DMZ pattern
but a federated research-network zone --- the same high-throughput
research-data path realised at national, mutualised scale ---
\emph{co-designed} with the \loc{nren-name} and shared across
partner research organisations, with a dedicated medical enclave for
biomedical workloads where required.

\mainref{the sections \nameref{subsec:lux-sdmz-status}, \nameref{sec:couts} and \nameref{sec:transition}, and the standalone Technical Reference~\cite{ng-tech-ref-2026} (vendor-neutral architecture, sovereignty grid, migration playbook)}

%============================================================================
%\newpage
\section{Current regulatory and operational context}
%============================================================================

\begin{enumerate}
  \item \textbf{AI as a structural tool for research and teaching:}
    research and teaching access to major cloud-based and open AI
    services (commercial APIs and open model hubs) is now part of
    the standard scientific working environment.
    Under the \loc{ai-regulation}'s risk-based framework --- the
    technology-governance obligation (F5) of the Regulatory Applicability annex
    (Part~\ref{part:reg-applicability} of the complete dossier) --- the
    resulting obligations
    depend on the concrete use case and the user's role; many research
    uses of commercial AI services fall outside the high-risk categories.
    Under the proposed cooperative researcher-centred model
    (PI data classification + institutional register maintained by
    the central IT function + periodic audit by the data-protection
    function), the cross-border-transfer obligation (F4) bites: in the worked
    EU instantiation, the \emph{Schrems~II} cumulative regime applies to
    personal-data transfers~\cite{schrems2020}.
    The AI regime's research-and-development exemption (where
    one exists --- e.g.\ Art.~2(6) of the EU AI Act) covers
    bespoke research tools, not the use of commercial APIs.
  \item \textbf{The applicable regulatory framework demands it:} the
    data-protection security-of-processing obligation (F1) requires
    \emph{proportionate} measures --- \gls{gdpr} Article~32 in the worked EU
    instantiation; under the cybersecurity-baseline regime (F2), instantiated
    here by the national transposition of \gls{nis2}~\cite{pl-8364}, the
    institution may be classified by the \loc{nis2-supervisory-authority}
    as an \textbf{essential entity}
    (holistic perimeter, proactive supervision, mandatory audits,
    annual maturity report via the \loc{national-notification-platform}~\cite{serima})~\cite{ilr-ul-essential};
    sector cybersecurity guidance from the relevant sector authority
    treats cybersecurity as a strategic governance
    risk~\cite{jisc-guidance}; a segmented architecture aligned with
    \gls{nis2} Article~21 proportionality is the appropriate response. The
    instruments named here are the EU instantiation of a
    jurisdiction-invariant obligation spine; an adopter reads its own
    regime's instruments for these obligations in the Regulatory
    Applicability annex (Part~\ref{part:reg-applicability} of the
    complete dossier).
  \item \textbf{Reference universities providing comparable openness:}
    reference research institutions
    provide their researchers with the autonomy and infrastructure
    described here. The \loc{nren-name} and its
    international backbone interconnect
    make the same architectural primitives accessible at the
    institution's national scale.
\end{enumerate}

\mainref{the sections \nameref{sec:reglementation}
and \nameref{sec:ia}}

%============================================================================
\section{A matter of institutional coherence}
%============================================================================

\begin{actionbox}[Call to action]
The university cannot proclaim academic freedom in its statutes and impose
a uniform digital lockdown in its daily practice. Providing an
\textbf{open, segmented and proportionately secured} framework means
honouring the university spirit while \textbf{strengthening security}
where it is truly necessary --- rather than weakening the institution
through blanket restrictions that drive talent away and fuel Shadow~IT.

\medskip
\textbf{Short-term objectives:}
\begin{enumerate}
  \item That the \textbf{unit-level deliberative body}, subsequently the
    \textbf{institution-wide council}, and finally
    \textbf{\loc{governing-board}} endorse a position statement derived
    from the present executive summary, \emph{``\suitpolicytitle''};
  \item That the \textbf{executive leadership}, in collaboration with IT
    experts and academic representatives, define a short- to
    medium-term implementation plan aligned with the architecture
    described in the accompanying technical report and developed
    in detail in the standalone Technical Reference~\cite{ng-tech-ref-2026}.
\end{enumerate}
\end{actionbox}

\vfill

\noindent\rule{\textwidth}{0.4pt}

{\footnotesize\textbf{Accompanying documents:} \emph{``Infrastructure and
Organisation of Information Systems for the institution ---
Analysis, State of the Art and Propositions''}. The report provides the full state of the art, an international set of
university case studies, European regulatory analysis, detailed
architecture, cost-benefit analysis, transition plan with risk register and
KPIs, and the complete legal and normative foundations for the \gls{icradp}.
A standalone Technical Reference~\cite{ng-tech-ref-2026} addressed to
chief technology officers and information-security officers develops
a five-dimension sovereignty grid, a vendor-neutral reference
architecture demonstrated on open foundations, and a wave-based
migration playbook. It is intended as the shared technical baseline
for the cross-institution federation anticipated in phase~2 of
this initiative.}

\smallskip

{\footnotesize\textbf{Note on the process:}
\textit{This document is an always living document.
It is provided for constructive feedback and revision in order to improve its quality.
The engineering of this document was partially supported by digital tools, some of which incorporate Artificial Intelligence features.}}
